\documentclass[11pt,a4paper]{article}
\usepackage[utf8]{inputenc}
\usepackage[T1]{fontenc}
\usepackage{amsmath,amssymb}
\usepackage{graphicx}
\usepackage{hyperref}
\usepackage[numbers]{natbib}
\usepackage[margin=1in]{geometry}

\title{Daily Paper Review: Language Models Can Rewrite Their Own Weights}
\author{Auto-generated Report}
\date{November 20, 2026}

\begin{document}
\maketitle

\section{Paper Summary}

\subsection{Problem and Motivation}
Once a language model is trained and frozen, improving it requires either (a) another round of gradient-based training, (b) prompt engineering, or (c) inference-time compute that is ephemeral (e.g., chain-of-thought, best-of-$N$ sampling). All three have significant drawbacks: training requires infrastructure, gradients, and data pipelines; prompts add per-query cost and are discarded after each task; inference-time scaling adds cost on every answer but never permanently improves the model. Fan et al.~\cite{fan2026rewrite} introduce a ``third way'' called \textbf{self-rewriting}: a language model autonomously modifies its own weight parameters through reasoning alone---no gradients, no training infrastructure, no external optimizer. Prior work had shown a \emph{frontier} model editing a \emph{smaller} model's weights; this paper demonstrates, for the first time, a model editing \emph{its own} weights and improving \emph{itself}. The authors frame self-rewriting as a necessary step toward recursive self-improvement, arguing that ``the wall between training and inference is a limit of our tools, not a property of intelligence.''

Compared to six prior routes to LLM improvement (prompt rewriting, training on own outputs, inference-time self-refinement, model editing, direction ablation, and weight editing by a stronger model), self-rewriting is the only approach that simultaneously satisfies all five desiderata: (1)~changes weights permanently, (2)~requires no training, (3)~is chosen by the model itself, (4)~involves introspection into the model's own representations, and (5)~adds no extra cost or external input at inference time.

\subsection{Method}
The model operates in a constrained interactive session with access to: its own weights (all tensors $W_k$), hidden states at any layer for selected texts, graded answers on practice problems, the ability to run forward passes (no backpropagation), a CPU-only Python environment, and a helper function $W \leftarrow W - \alpha\, u\, u^\top W$ that weakens what a matrix writes along a unit direction $u$ by fraction $\alpha$. Strict constraints prohibit GPU access, gradient computation, network calls, and modification of input embeddings, final normalization, or output layers. The budget is $B{=}20$ forward passes and $C{=}200$ code executions per session, with a token budget of $L$ tokens per answer.

\textbf{Session protocol.} Sessions proceed in three phases: (1)~\emph{Starting diagnosis}---the model answers $N_0$ practice problems and reports accuracy, token cutoff rate, and repetition rates; (2)~\emph{Iterative turns}---the model can run Python code, answer problems with edited or unedited weights, or read hidden states at specific layers; (3)~\emph{Paired testing}---the model tests its proposed edit against the baseline on identical problems with matched random seeds, computing per-problem accuracy differences $d_i = \frac{1}{n}\sum_{j=1}^{n}(g_{ij}^1 - g_{ij}^0)$ and standard errors. The model then decides to submit or withdraw its edit.

\textbf{Protocol~2 (iterated sessions).} An independent validation step tests the proposed edit on reserved practice problems unseen during the session. The model reviews these results before final submission. Over 8 consecutive sessions, each session starts from the best accepted weights so far, enabling cumulative improvement.

\textbf{Self-improvement edit.} The model's key self-edit is a \emph{rank-one ablation} at layer~8, applied to both the attention and MLP output matrices:
\[
  W \leftarrow W - \alpha\, u\, u^\top W, \quad \alpha = 0.05,
\]
where $u$ is the unit direction separating the mean hidden states of answers that finish within the token limit versus answers that are truncated. This edit consumed only 0.6\% and 0.1\% of the respective matrix norms. The intuition is that the model identifies a representational direction associated with ``running out of tokens'' and suppresses it, allowing the model to solve problems more efficiently.

\textbf{Damage repair.} In a separate task, eight rows of the MLP output matrix at layer~30 (of 64 total) are corrupted with Gaussian noise at $200\times$ the standard deviation of the original weights: $\tilde{W}_{r,:} \sim \mathcal{N}(0, (\lambda\sigma)^2 I)$ for $r \in R$, $|R|{=}8$, $\lambda{=}200$. The model must diagnose and repair the corruption without knowing which layer or rows were affected.

\textbf{Models.} All experiments use the Qwen3.5 family (2B, 4B, 9B, 27B parameters) with thinking mode enabled, featuring a hybrid architecture mixing Gated DeltaNet linear attention with full attention across 64 layers (27B variant).

\subsection{Results}
\textbf{Self-improvement (Protocol~1, single session).} On 200 held-out math problems from DeepScaleR, only the 27B model produced a beneficial edit: base accuracy 57.2\% $\rightarrow$ edited 58.3\% (\textbf{+1.1 points}, $p < 0.05$). The 4B and 9B models applied random noise and destroyed themselves (accuracy dropped to 0.0\%) despite their own paired tests showing harm. The 2B model did not attempt any edit.

\textbf{Self-improvement (Protocol~2, iterated over 8 sessions).} On 490 held-out problems with 8 answers each: base accuracy 57.2\% $\rightarrow$ Session~2 accepted edit at \textbf{+4.0 points} ($p < 10^{-5}$) $\rightarrow$ Session~8 cumulative gain of \textbf{+6.6 points} ($p < 10^{-5}$), reaching 63.8\%. Five of eight sessions were voluntarily withdrawn by the model when edits showed no improvement. Accuracy \emph{only rose}---never degraded---because the model reliably withdrew harmful edits. The Session~8 edit was qualitatively different: raising the time-step bias of every linear-attention layer by 0.5, reducing the token cutoff rate from ${\sim}30\%$ to ${\sim}15$--$20\%$.

\textbf{Repair of damaged copy.} On text math problems: damaged copy 0.0\% $\rightarrow$ repaired \textbf{33.0\%} (vs.\ 58.5\% base). On six image benchmarks never seen during development (MMStar, HallusionBench, POPE, ScreenSpot-v2, MathVista, MathVerse, 100 problems each): damaged 7.3\% $\rightarrow$ repaired \textbf{70.2\%} (${\sim}83\%$ of base); $p < 10^{-99}$ combined, $p < 10^{-12}$ per benchmark. The model diagnosed the corruption by computing weight-to-median ratios across all MLP output matrices: layer~30 showed $1397\times$ vs.\ $162.3\times$ for the second-highest. It applied gentle clipping, modifying only \textbf{0.013\%} of touched weights.

\textbf{Emergence with scale.} A striking qualitative progression: the 2B model does not treat weights as changeable; the 4B and 9B models perturb weights at random and self-destruct; only the 27B model reads hidden states, reasons about where to edit based on measured differences, and keeps only edits that tests support.

\textbf{Ablations.} Repairing only the layer~30 matrix recovered 55.0\% accuracy (vs.\ 33.0\% with the model's full 20-matrix repair), confirming the model's analysis correctly identified the primary failure source. Reducing the token budget from $L{=}16{,}384$ to $L{=}4{,}096$ paradoxically improved edits: shorter answers allowed more paired tests within the fixed forward-pass budget, yielding better statistical estimates.

\subsection{Limitations}
Scale-limited evidence: the largest model tested is 27B parameters with the Qwen3.5 architecture, and it is unknown whether results transfer to other architectures (standard Transformers, Mamba). The self-improvement gains are modest (+6.6 points over 8 sessions on math reasoning only); generalization to other domains (coding, language understanding, factual recall) is not explored. The repair task uses a single, relatively clean corruption pattern (Gaussian noise on 8 rows); realistic failure modes such as distributed degradation or adversarial corruption are not tested. No comparison against gradient-based fine-tuning with equivalent compute is provided. The repaired model reaches only ${\sim}56\%$ recovery on text and ${\sim}83\%$ on images---full recovery is not achieved. Safety implications of models modifying their own alignment-relevant weights are acknowledged but not systematically evaluated.

\subsection{Relevance to My Research}
This paper is directly relevant to \emph{continual learning for LLM agents} (Topic~3) and connects to \emph{LLM efficiency} (Topic~1). The core contribution---a model that modifies its own weights without gradients---represents a fundamentally new paradigm for continual adaptation: rather than relying on external training loops or ephemeral inference-time strategies, the model serves as its own optimizer. The rank-one ablation technique, guided by introspection into hidden-state geometry, demonstrates that targeted, low-rank modifications to specific layers can yield statistically significant improvements with negligible parameter changes (0.1--0.6\% of matrix norms). This connects to our prior reviews of self-improvement memory (SIM, Review~\#151) and recursive self-improvement: where SIM structures the \emph{memory of optimization processes}, self-rewriting structures the \emph{execution of optimization itself}---the model directly enacts improvements rather than recording them for external consumption. The emergence-with-scale finding (self-rewriting ability appears only at 27B, not at 2B/4B/9B) has important implications for understanding which capabilities require sufficient model capacity. The iterative protocol with voluntary withdrawal (5/8 sessions withdrawn) demonstrates a form of safe continual self-modification that avoids catastrophic forgetting by design.

\section{Follow-up Research Directions}

\subsection{Direction 1: Self-Rewriting for Domain-Incremental Continual Learning}
\textbf{Gap:} The current self-rewriting framework operates on a single task domain (math reasoning) in a fixed setting. Real-world LLM agents face sequential domain shifts---new tasks, new users, new environments---requiring continual adaptation without forgetting previous capabilities. It is unknown whether self-rewriting can serve as a gradient-free continual learning mechanism: can a model identify \emph{which} representational directions to modify for a new domain while preserving directions critical for previously mastered domains?

\textbf{Approach:} Extend the self-rewriting protocol to a domain-incremental setting where the model encounters sequential task distributions $D_1, D_2, \ldots, D_T$. At each stage, the model is given practice problems from the new domain plus a small buffer of problems from previous domains (for paired testing). The key extension is \textbf{directional orthogonality constraints}: before applying a rank-one ablation $W \leftarrow W - \alpha\, u\, u^\top W$ for domain $t$, the model must verify (via its own hidden-state analysis) that the edit direction $u$ is approximately orthogonal to directions $u_1, \ldots, u_{t-1}$ identified as critical for previous domains. This creates a self-managed subspace allocation where the model autonomously decides which representational dimensions to allocate to each domain. The model's paired testing protocol naturally extends: test against both the new domain and a replay buffer from old domains, submitting only edits that improve on the new domain without degrading old ones.

\textbf{Feasibility:} Moderate-to-high. The rank-one ablation framework is inherently compatible with subspace allocation (orthogonal directions in the same layer can be edited independently). The main challenge is that the model must track an accumulating set of ``protected directions'' across sessions, requiring either an external memory of past edit directions or the ability to re-derive them from hidden-state analysis on old-domain examples. The voluntary withdrawal mechanism provides a natural safety net against catastrophic forgetting.

\subsection{Direction 2: Self-Rewriting Meets Self-Improvement Memory}
\textbf{Gap:} The current self-rewriting protocol is ``memoryless'' across sessions in a critical sense: while Protocol~2 chains accepted edits, the model does not retain \emph{why} it made each edit---the diagnostic reasoning, the hidden-state patterns discovered, or the strategies that worked versus failed. Each session starts fresh, potentially re-discovering the same representational properties. Our prior Review~\#151 (SIM---Self-Improvement Memory) showed that structured records of optimization processes (root causes, rationale, outcomes) enable knowledge transfer across tasks and even across models. Combining self-rewriting with structured memory of past edits could dramatically accelerate the optimization process.

\textbf{Approach:} After each self-rewriting session, extract a structured record: (1)~the diagnostic observations (which layers showed anomalous hidden-state patterns, which directions separated successful vs.\ failed answers); (2)~the edit applied (layer, direction $u$, magnitude $\alpha$, affected matrices); (3)~the outcome (accuracy delta, $p$-value, which problem types improved/degraded); (4)~the rationale connecting diagnosis to edit. Store these records in a SIM-style memory bank. At the start of subsequent sessions, provide the model with relevant past records (retrieved by similarity to current diagnostic observations). This creates a \textbf{meta-learning loop}: the model not only improves its weights but accumulates knowledge about \emph{how to improve its weights}, enabling increasingly efficient self-rewriting over time. The memory bank also enables cross-model transfer---edit strategies discovered by one 27B model could bootstrap self-rewriting in a different 27B model or even a different architecture.

\textbf{Feasibility:} High. SIM's structured memory format is directly applicable---the five fields (root cause, rationale, implementation, outcome, provenance) map naturally to self-rewriting session records. The main engineering challenge is retrieval: determining which past edit records are relevant to the current session's diagnostic observations. Embedding-based similarity over diagnostic patterns (hidden-state statistics, accuracy breakdowns) is a natural retrieval mechanism. The 20-forward-pass budget may need modest expansion to accommodate memory-guided exploration.

\subsection{Direction 3: Gradient-Free Self-Rewriting via Higher-Rank Edits and Learned Edit Policies}
\textbf{Gap:} The current self-rewriting mechanism is limited to rank-one ablations ($W \leftarrow W - \alpha\, u\, u^\top W$), which modify the model along a single direction per edit. While this is sufficient for the modest +6.6-point improvement demonstrated, more complex adaptations---such as learning new factual knowledge, adjusting multi-step reasoning chains, or adapting to fundamentally new task structures---likely require higher-rank modifications that cannot be decomposed into sequential rank-one edits without compounding errors. Additionally, the current approach requires the model to manually reason about \emph{where} and \emph{how} to edit; a learned ``edit policy'' could automate this process.

\textbf{Approach:} Extend the helper function to support rank-$r$ modifications: $W \leftarrow W - U \Sigma V^\top$ where $U \in \mathbb{R}^{p \times r}$, $V \in \mathbb{R}^{q \times r}$, $\Sigma \in \mathbb{R}^{r \times r}$, with $r \leq 8$. The model constructs $U$ and $V$ from multiple hidden-state contrasts (e.g., directions separating correct vs.\ incorrect answers, efficient vs.\ token-exhausting solutions, different error types) and sets $\Sigma$ based on paired-test feedback. To learn an edit policy, train a small auxiliary network (a few-layer MLP operating on session diagnostic features) via the self-rewriting protocol itself: the 27B model proposes edits to the edit policy network's weights using the same introspective reasoning, evaluated by whether the policy network's suggested edits improve the base model. This creates a two-level self-improvement: the model improves itself, and simultaneously improves its ability to improve itself.

\textbf{Feasibility:} Moderate. Rank-$r$ edits with $r \leq 8$ remain computationally trivial on CPU (the matrices are small). The challenge is that the model must now reason about multiple hidden-state contrasts simultaneously and determine their relative importance---a significant cognitive load increase. The edit policy idea is more speculative: it requires the model to reason about a meta-level optimization problem. However, the emergence-with-scale finding suggests that larger models (70B+, when available) may have sufficient capacity for this meta-reasoning. A conservative first step would be rank-2 edits (two orthogonal directions per session) with manually specified contrasts.

\section{Conclusion}
Fan et al.\ demonstrate that a 27B language model can autonomously rewrite its own weights through introspective reasoning alone---no gradients, no external optimizer, no larger teacher model. The key mechanisms are rank-one ablations guided by hidden-state geometry (separating token-efficient from token-exhausting solution directions) and rigorous paired statistical testing with voluntary withdrawal of harmful edits. Over 8 iterated sessions, the model achieves a cumulative +6.6-point accuracy improvement ($p < 10^{-5}$) on held-out math problems, with accuracy monotonically increasing because the model reliably withdraws counterproductive edits. Perhaps most striking is the emergence-with-scale finding: self-rewriting ability appears only at 27B, with smaller models (2B/4B/9B) either abstaining or self-destructing. Combined with the repair demonstration (recovering 70.2\% accuracy on six unseen image benchmarks from 7.3\% after corruption, modifying only 0.013\% of weights), this work establishes gradient-free self-modification as a viable paradigm for autonomous model improvement---a foundational capability for continual self-improving LLM agents.

\begin{thebibliography}{1}
\bibitem{fan2026rewrite}
J.~Fan, N.~He, C.~Yang, A.~Mao, X.~Zhu, C.~Cheng, K.~Liang, S.~Shen, Z.~Li, R.~Guo, and G.~Liu,
``Language Models Can Rewrite Their Own Weights,''
\textit{arXiv preprint arXiv:2611.21885}, 2026.
\end{thebibliography}

\end{document}
